\section{Monte Carlo 与随机逼近把回报变成增量更新}
\label{sec:16-Control-and-Reinforcement-Learning/02-Reinforcement-Learning/03}

你接手一套配送调度系统，想评估现行的派单策略值多少钱。上一节的方法在这里全部作废：没人写得出需求分布和路况转移核，那个对所有 \(s'\) 求和的式子根本无从落笔。能做的只有让策略真跑一天，晚上把当天的折扣回报记下来。

于是问题换了个形状。一天只能拿到一个数字，而这个数字里塞满了天气、订单波动、某个骑手迟到的噪声。你要多少天才能说清这个策略的价值？每天来一个新数字，旧的估计该往新数字挪多少？挪多了追着噪声跑，挪少了几个月都反应不过来。这一节要说的是：一旦把期望换成样本，收敛条件就从``压缩系数小于一''换成了一对关于步长的级数条件。

\subsection{一幕的回报只是价值的一个样本}

固定策略 \(\pi\)，从时刻 \(t\) 起的折扣回报仍是

\[
G_t=R_{t+1}+\gamma R_{t+2}+\gamma^2R_{t+3}+\cdots,
\]

而状态价值按定义就是它的条件期望 \(V^\pi(s)=\E_\pi[G_t\given S_t=s]\)。所以每次访问 \(s\) 之后观察到的那个 \(G_t\)，是 \(V^\pi(s)\) 的一个无偏样本。收集 \(N(s)\) 个这样的回报，取平均

\[
\widehat V(s)=\frac1{N(s)}\sum_{i=1}^{N(s)}G^{(i)},
\]

由大数定律它会收敛到 \(V^\pi(s)\)。这条路径与\hyperref[sec:05-Probability/08-Limit-Theorems/05]{``Monte Carlo：从随机样本到数值答案''}里用抽样算积分是同一件事，只是被积的东西现在是一整条轨迹的回报。

首次访问法只用状态在一幕中第一次出现之后的回报，每次访问法则把每次出现都算上。两者在合适的遍历条件下都相合，但有限样本下的相关结构和方差不一样——一幕之内多次访问同一状态得到的回报是高度相关的，把它们当独立样本去估方差会明显偏小。

Monte Carlo 的好处很干净：不需要转移概率，目标里也不含当前的价值估计，所以没有自举误差。麻烦同样干净：必须等到一幕结束，而 \(G_t\) 把沿途每一步的奖励噪声都累加了进去。时域越长、\(\gamma\) 越接近 1，这个方差越吓人。配送系统里一天有几百次派单，全部噪声压进一个数字，靠一天的观测什么也说明不了。

\subsection{样本平均可以写成一行增量更新}

不必攒够 \(N\) 个回报再平均。样本均值本身满足递推

\[
\bar G_n
=\bar G_{n-1}+\frac1n\bigl(G_n-\bar G_{n-1}\bigr).
\]

括号里是新观测与旧估计的落差，\(1/n\) 是步长。这个形状会在后面反复出现：

\[
\text{新估计}
=\text{旧估计}+\alpha_t\times(\text{目标}-\text{旧估计}).
\]

内存从 \(O(N)\) 降到 \(O(1)\)，这只是顺带的好处。要紧的是把 \(1/n\) 换成常数 \(\alpha\) 之后会发生什么：旧数据的权重按 \((1-\alpha)^k\) 指数衰减，估计变成了一个滑动加权平均。它能追踪缓慢漂移的环境，但不再收敛到任何固定值，而是在目标附近持续抖动。

同一行更新式，两种截然不同的行为。这件事不能靠看更新式的外形判断，只能看步长序列。

\subsection{步长决定你估的是固定期望还是移动目标}

\begin{center}
\begin{tikzpicture}[>=stealth, scale=0.9]
% 左：两个级数的走向
\draw[->] (0,0) -- (5.6,0);
\draw[->] (0,0) -- (0,3.3);
\node[anchor=north] at (2.8,-0.15) {\small \(t\)};
\draw[domain=0.02:5,samples=80,smooth] plot ({\x},{0.95*ln(1+3.4*\x)});
\draw[dashed,domain=0:5,samples=80,smooth] plot ({\x},{1.45*(1-exp(-2.2*\x))});
\node[anchor=west] at (2.5,2.85) {\small \(\sum\alpha_t\to\infty\)};
\node[anchor=west] at (2.5,1.15) {\small \(\sum\alpha_t^2<\infty\)};
% 右：两种步长下的估计轨迹
\begin{scope}[shift={(6.8,0)}]
\draw[->] (0,0) -- (5.6,0);
\draw[->] (0,0) -- (0,3.3);
\node[anchor=north] at (2.8,-0.15) {\small \(t\)};
\draw[dotted] (0,2.0) -- (5.3,2.0);
\node[anchor=west] at (5.3,2.0) {};
\draw[smooth] plot coordinates {(0,0.4) (0.4,1.5) (0.8,2.6) (1.2,1.7) (1.6,2.35) (2.0,1.82) (2.4,2.15) (2.8,1.92) (3.2,2.08) (3.6,1.96) (4.0,2.03) (4.4,1.99) (4.8,2.01) (5.2,2.0)};
\draw[dashed,smooth] plot coordinates {(0,0.4) (0.4,1.6) (0.8,2.75) (1.2,1.55) (1.6,2.55) (2.0,1.62) (2.4,2.48) (2.8,1.68) (3.2,2.44) (3.6,1.70) (4.0,2.40) (4.4,1.74) (4.8,2.36) (5.2,1.78)};
\node[anchor=west] at (2.2,1.05) {\small 实线 \(\alpha_t=1/t\)};
\node[anchor=west] at (2.2,0.50) {\small 虚线 \(\alpha_t\equiv\alpha\)};
\end{scope}
\end{tikzpicture}
\end{center}

\emph{左：累积步长必须发散、累积平方必须收敛；右：满足这一对条件的步长把估计钉在目标上，固定步长只把它围在目标附近。}

图的右半边就是配送系统里两种运维选择的写照。实线那条最终停住，三个月后它对新开的仓库却几乎无感——步长已经小到推不动了。虚线那条永远跟得上变化，也就永远在抖，你没法用它的读数去比较两个策略的优劣，因为差异可能全在抖动幅度之内。

左半边解释了为什么必须是这一对条件。第一个级数 \(\sum_t\alpha_t=\infty\) 要求总步长发散：只要还有偏差没消掉，迭代就得留有把自己挪过去的余力；步长衰减太快，级数收敛，估计会在半路冻住，停在一个与初值有关的错误位置。第二个级数 \(\sum_t\alpha_t^2<\infty\) 要求噪声注入的累计能量有限：每步噪声以 \(\alpha_t\) 的比例进来，方差按 \(\alpha_t^2\) 累加，平方和有限才能保证抖动最终被压平。两者拉扯的是同一个 \(\alpha_t\)：\(\alpha_t=1/t\) 恰好卡在中间，\(\alpha_t=1/\sqrt{t}\) 满足第一条而违反第二条，常数步长两条都不满足第二条。判断一个级数落在哪一边，用的就是\hyperref[sec:02-Calculus/05-Sequences-and-Series/03]{``正项级数的比较''}里那套标准。

\subsection{Robbins--Monro 条件不是完整的收敛定理}

一般的随机逼近要解 \(h(\theta)=\E[H(\theta,\xi)]=0\)，但每次只观察到带噪的 \(H(\theta_t,\xi_{t+1})\)，更新写成

\[
\theta_{t+1}=\theta_t-\alpha_tH(\theta_t,\xi_{t+1}).
\]

上面那对级数条件就是 Robbins--Monro 条件。要看清它为什么管用，取一个最简特例：

\[
Z_{t+1}=(1-\alpha_t)Z_t+\alpha_t\xi_{t+1},
\qquad 0<\alpha_t\le 1,
\]

其中噪声相对于截至 \(t\) 的信息满足 \(\E[\xi_{t+1}\given\mathcal F_t]=0\) 与 \(\E[\xi_{t+1}^2\given\mathcal F_t]\le C\)。此时在两个级数条件下 \(Z_t\to 0\) 几乎必然成立。骨架分两半：旧误差每步被乘以 \(1-\alpha_t\)，连乘的对数是 \(-\sum_t\alpha_t\) 的量级，第一个条件让它趋于 \(-\infty\)，所以初始误差被彻底抹掉；新噪声只以 \(\alpha_t\) 的权重进入，其累积方差由 \(\sum_t\alpha_t^2\) 控住，第二个条件让它有限。收缩不停、噪声可和，两股力量各管一头。

条件对账要老实做。这个特例假设了收缩中心固定、噪声条件均值为零、条件方差一致有界。真实的价值更新一条都不自动满足：收缩中心是 Bellman 算子的像，会随迭代移动；样本由 Markov 链产生，前后相关，条件均值为零要靠额外论证；迭代还必须先证明有界，否则一切估计都是空的。所以标准做法是把这些写成完整定理的前提——平均动力系统有稳定根、迭代几乎必然有界、噪声可测且矩条件成立——再套 Dvoretzky 型结论。更新式长得像加权平均，不等于收敛结论可以直接搬。

\hyperref[sec:09-Convex-Optimization/08-Nonconvex-and-Stochastic-Optimization/04]{``SGD 在期望中下降但噪声会留下平台''}从光滑目标的下降引理出发，把固定步长留下的那圈驻点邻域算成了方差与步长的显式函数。这里的情形更麻烦一层：那里的目标函数不动，这里的目标随着估计一起移动。

\subsection{要改进策略就必须探索，而真实系统会限制怎么探索}

只评估一个固定策略时，样本从那个策略里来就够了。一旦要改进策略，就得知道每个状态下\emph{各个}动作的价值：

\[
Q^\pi(s,a)=\E_\pi[G_t\given S_t=s,A_t=a].
\]

从没被选过的动作没有任何回报数据，它的估计只能停在初值上，而贪心又只会挑估计高的动作——初值偏低的好动作可能永远得不到翻身的机会。Exploring starts 假设每个状态动作对都有机会当一幕的起点，理论上方便，实际中很少可行。更常用的是 \(\varepsilon\)-greedy：当前最优动作拿 \(1-\varepsilon\) 的概率，剩下的 \(\varepsilon\) 均分给其他动作。

探索保证``能看到''，不保证``及时看到''。稀有状态、延迟奖励和高风险动作都会让均匀随机的探索付出巨大的样本代价。\(\varepsilon\) 降到零太快，可能永久漏掉更好的动作；永不下降，则要一直付出执行次优动作的成本。

在真实冷库里，\(\varepsilon\)-greedy 的``随机试一下''可能就是突然停掉制冷，或者让压缩机满功率运行。一次危险探索会毁掉一批药品，频繁切换会磨损设备，而真实系统没有复位键，也不能为了覆盖极端状态去主动制造一次高温事故。可行的做法是把探索挪到别处：先在经过验证的动力学模型或数字孪生里试，用离线数据在历史策略的支撑范围内做评估，给动作和温度设硬约束，再留一个能随时接管的安全控制器。这些手段限制的是允许收集哪一类数据，它们不能被理解成给同一个无约束算法换个名字。

\subsection{Monte Carlo 与 TD 交换的是偏差和方差}

Monte Carlo 的目标是完整回报 \(G_t\)，无偏但方差大，而且要等一幕结束。时序差分换成一步目标

\[
R_{t+1}+\gamma V(S_{t+1}),
\]

不必等待，方差通常小得多，而目标里从此含有当前的估计 \(V\)——这就是自举。估计的误差会顺着目标进入下一次更新。

这不是``准确解''与``近似解''的区别。表格情形下，TD 在适当条件下收敛到与 Monte Carlo 相同的 Bellman 不动点，两者只是走向同一答案的不同路线。差别要到函数逼近、离策略数据与自举三样凑齐时才会变得危险。下一节先把 TD 的思路推到控制上，看看用采样替换期望之后，那个留在方程里的 \(\max\) 会带来什么。
